\scanpage{23}
% archive foliation: 22
\begin{center}- 87 -\end{center}
n° 268
\[
Z_{(-x,u)}\cdot(0\times A)=0\times Z(x)_u.
\]
On a une formule semblable en prenant $w=-y$. Cela montre que
\[
X_u-X=\operatorname{pr}_A[T\cdot(0\times A)]
\]
Mais d'après le théorème 4, Corollaire 1, on a $T\sim0$ sur $J\times A$, et par suite $X_u-X\sim0$ sur $A$ d'après F-VIII$_2$ Th. 4, Coro. 1. Comme $u$ est arbitraire, cela démontre notre proposition.

3-5 - \uline{Le Noyau du théorème du Carré.}

Soit $V^r$ une variété complète, sans singularités de codimension 1, définie sur $k$. Soit $P$ un point simple de $V$, rationnel sur $k$. Soit $\Omega$ le domaine universel, et soient $t_1,\ldots,t_r$ des fonctions sur $V$, définies sur $k$, qui soient des paramètres uniformisant de $P$ sur $V$. Désignons par $\unsure{\mathfrak{o}}$ l'anneau local de $P$ sur $\Omega$, et $\mathfrak{m}$ son idéal maximal.

Alors $t_1,\ldots,t_r$ sont des générateurs de $\mathfrak{m}/\mathfrak{m}^2$ sur $\mathfrak{o}/\mathfrak{m}$, et toute fonction $f$ s'exprime par une série de puissance en monômes des $t_i$ avec coefficients dans $\Omega$. On notera par $M_j(t)$ tous les monômes $t_1^{\mu_1}\ldots t_r^{\mu_r}$ da,s n'importe quel ordre $j=0,1,2,\ldots$ assujetti seulement à la condition que $\deg M_j\leq\deg M_{j+1}$. Nous pouvons alors écrire $f\in\mathfrak{o}$,
\[
f=\sum_{j=\infty}^{\infty}c_jM_j(t)\qquad c_j\in\Omega.
\]
Si $f$ est définie sur un corps $K\supset k$, alors les $c_j$ sont dans $K$. Le sous-anneau de $\mathfrak{o}$ des fonctions $f$ définies sur $K$ sera noté $\mathfrak{o}^K$ et son idéal maximal par $\mathfrak{m}^K$.

Le sous ensemble des fonctions de $\mathfrak{o}$ dont la série de puissance est telle que
\[
c_0=c_1=\cdots=c_N=0
\]
est évidemment un idéal $I_N$ de $\mathfrak{o}$, et on obtient de cette façon une suite décroissante d'idéaux dont l'intersection est réduite à zéro.
